%%% Hlavní soubor. Zde se definují základní parametry a odkazuje se na ostatní části. %%%

%% Verze pro jednostranný tisk:
% Okraje: levý 40mm, pravý 25mm, horní a dolní 25mm
% (ale pozor, LaTeX si sám přidává 1in)
\documentclass[12pt,a4paper]{report}
\setlength\textwidth{145mm}
\setlength\textheight{247mm}
\setlength\oddsidemargin{15mm}
\setlength\evensidemargin{15mm}
\setlength\topmargin{0mm}
\setlength\headsep{0mm}
\setlength\headheight{0mm}
% \openright zařídí, aby následující text začínal na pravé straně knihy
\let\openright=\clearpage

%% Pokud tiskneme oboustranně:
% \documentclass[12pt,a4paper,twoside,openright]{report}
% \setlength\textwidth{145mm}
% \setlength\textheight{247mm}
% \setlength\oddsidemargin{14.2mm}
% \setlength\evensidemargin{0mm}
% \setlength\topmargin{0mm}
% \setlength\headsep{0mm}
% \setlength\headheight{0mm}
% \let\openright=\cleardoublepage

%% Vytváříme PDF/A-2u
\usepackage[a-2u]{pdfx}

%% Přepneme na českou sazbu a fonty Latin Modern
\usepackage[czech]{babel}
\usepackage{lmodern}
\usepackage[T1]{fontenc}
\usepackage{textcomp}

%% Použité kódování znaků: obvykle latin2, cp1250 nebo utf8:
\usepackage[utf8]{inputenc}

%%% Další užitečné balíčky (jsou součástí běžných distribucí LaTeXu)
\usepackage{amsmath}        % rozšíření pro sazbu matematiky
\usepackage{amsfonts}       % matematické fonty
\usepackage{amsthm}         % sazba vět, definic apod.
\usepackage{bbding}         % balíček s nejrůznějšími symboly
			    % (čtverečky, hvězdičky, tužtičky, nůžtičky, ...)
\usepackage{bm}             % tučné symboly (příkaz \bm)
\usepackage{graphicx}       % vkládání obrázků
\usepackage{fancyvrb}       % vylepšené prostředí pro strojové písmo
\usepackage{indentfirst}    % zavede odsazení 1. odstavce kapitoly
\usepackage{natbib}         % zajištuje možnost odkazovat na literaturu
			    % stylem AUTOR (ROK), resp. AUTOR [ČÍSLO]
\usepackage[nottoc]{tocbibind} % zajistí přidání seznamu literatury,
                            % obrázků a tabulek do obsahu
\usepackage{icomma}         % inteligetní čárka v matematickém módu
\usepackage{dcolumn}        % lepší zarovnání sloupců v tabulkách
\usepackage{booktabs}       % lepší vodorovné linky v tabulkách
\usepackage{paralist}       % lepší enumerate a itemize
\usepackage{xcolor}         % barevná sazba

%%% Údaje o práci

% Název práce v jazyce práce (přesně podle zadání)
\def\NazevPrace{Evoluce robotů v simulovaném fyzikálním prostředí}

% Název práce v angličtině
\def\NazevPraceEN{Evolution of robots in a simulated physical environment}

% Jméno autora
\def\AutorPrace{Marek Bečvář}

% Rok odevzdání
\def\RokOdevzdani{2023}

% Název katedry nebo ústavu, kde byla práce oficiálně zadána
% (dle Organizační struktury MFF UK, případně plný název pracoviště mimo MFF)
\def\Katedra{Katedra softwaru a výuky informatiky}
\def\KatedraEN{Department of software and computer science education}

% Jedná se o katedru (department) nebo o ústav (institute)?
\def\TypPracoviste{Katedra}
\def\TypPracovisteEN{Department}

% Vedoucí práce: Jméno a příjmení s~tituly
\def\Vedouci{RNDr. František Mráz, CSc.}

% Pracoviště vedoucího (opět dle Organizační struktury MFF)
\def\KatedraVedouciho{Katedra softwaru a výuky informatiky}
\def\KatedraVedoucihoEN{Department of software and computer science education}

% Studijní program a obor
\def\StudijniProgram{Informatika}
\def\StudijniObor{Informatika se specializací Umělá inteligence}

% Nepovinné poděkování (vedoucímu práce, konzultantovi, tomu, kdo
% zapůjčil software, literaturu apod.)
\def\Podekovani{%
Poděkování.
}

% Abstrakt (doporučený rozsah cca 80-200 slov; nejedná se o zadání práce)
\def\Abstrakt{%
Abstrakt.
}
\def\AbstraktEN{%
Abstract.
}

% 3 až 5 klíčových slov (doporučeno), každé uzavřeno ve složených závorkách
\def\KlicovaSlova{%
{klíčová} {slova}
}
\def\KlicovaSlovaEN{%
{key} {words}
}

%% Balíček hyperref, kterým jdou vyrábět klikací odkazy v PDF,
%% ale hlavně ho používáme k uložení metadat do PDF (včetně obsahu).
%% Většinu nastavítek přednastaví balíček pdfx.
\hypersetup{unicode}
\hypersetup{breaklinks=true}

%% Definice různých užitečných maker (viz popis uvnitř souboru)
%%% Tento soubor obsahuje definice různých užitečných maker a prostředí %%%
%%% Další makra připisujte sem, ať nepřekáží v ostatních souborech.     %%%

%%% Drobné úpravy stylu

% Tato makra přesvědčují mírně ošklivým trikem LaTeX, aby hlavičky kapitol
% sázel příčetněji a nevynechával nad nimi spoustu místa. Směle ignorujte.
\makeatletter
\def\@makechapterhead#1{
  {\parindent \z@ \raggedright \normalfont
   \Huge\bfseries \thechapter. #1
   \par\nobreak
   \vskip 20\p@
}}
\def\@makeschapterhead#1{
  {\parindent \z@ \raggedright \normalfont
   \Huge\bfseries #1
   \par\nobreak
   \vskip 20\p@
}}
\makeatother

% Toto makro definuje kapitolu, která není očíslovaná, ale je uvedena v obsahu.
\def\chapwithtoc#1{
\chapter*{#1}
\addcontentsline{toc}{chapter}{#1}
}

% Trochu volnější nastavení dělení slov, než je default.
\lefthyphenmin=2
\righthyphenmin=2

% Zapne černé "slimáky" na koncích řádků, které přetekly, abychom si
% jich lépe všimli.
\overfullrule=1mm

%%% Makra pro definice, věty, tvrzení, příklady, ... (vyžaduje baliček amsthm)

\theoremstyle{plain}
\newtheorem{veta}{Věta}
\newtheorem{lemma}[veta]{Lemma}
\newtheorem{tvrz}[veta]{Tvrzení}

\theoremstyle{plain}
\newtheorem{definice}{Definice}

\theoremstyle{remark}
\newtheorem*{dusl}{Důsledek}
\newtheorem*{pozn}{Poznámka}
\newtheorem*{prikl}{Příklad}

%%% Prostředí pro důkazy

\newenvironment{dukaz}{
  \par\medskip\noindent
  \textit{Důkaz}.
}{
\newline
\rightline{$\qedsymbol$}
}

%%% Prostředí pro sazbu kódu, případně vstupu/výstupu počítačových
%%% programů. (Vyžaduje balíček fancyvrb -- fancy verbatim.)

\DefineVerbatimEnvironment{code}{Verbatim}{fontsize=\small, frame=single}

%%% Prostor reálných, resp. přirozených čísel
\newcommand{\R}{\mathbb{R}}
\newcommand{\N}{\mathbb{N}}

%%% Užitečné operátory pro statistiku a pravděpodobnost
\DeclareMathOperator{\pr}{\textsf{P}}
\DeclareMathOperator{\E}{\textsf{E}\,}
\DeclareMathOperator{\var}{\textrm{var}}
\DeclareMathOperator{\sd}{\textrm{sd}}

%%% Příkaz pro transpozici vektoru/matice
\newcommand{\T}[1]{#1^\top}

%%% Vychytávky pro matematiku
\newcommand{\goto}{\rightarrow}
\newcommand{\gotop}{\stackrel{P}{\longrightarrow}}
\newcommand{\maon}[1]{o(n^{#1})}
\newcommand{\abs}[1]{\left|{#1}\right|}
\newcommand{\dint}{\int_0^\tau\!\!\int_0^\tau}
\newcommand{\isqr}[1]{\frac{1}{\sqrt{#1}}}

%%% Vychytávky pro tabulky
\newcommand{\pulrad}[1]{\raisebox{1.5ex}[0pt]{#1}}
\newcommand{\mc}[1]{\multicolumn{1}{c}{#1}}


%% Titulní strana a různé povinné informační strany
\begin{document}
%%% Titulní strana práce a další povinné informační strany

%%% Titulní strana práce

\pagestyle{empty}
\hypersetup{pageanchor=false}

\begin{center}

\centerline{\mbox{\includegraphics[width=166mm]{../img/logo-cs.pdf}}}

\vspace{-8mm}
\vfill

{\bf\Large BAKALÁŘSKÁ PRÁCE}

\vfill

{\LARGE\AutorPrace}

\vspace{15mm}

{\LARGE\bfseries\NazevPrace}

\vfill

\Katedra

\vfill

{
\centerline{\vbox{\halign{\hbox to 0.45\hsize{\hfil #}&\hskip 0.5em\parbox[t]{0.45\hsize}{\raggedright #}\cr
Vedoucí bakalářské práce:&\Vedouci \cr
\noalign{\vspace{2mm}}
Studijní program:&\StudijniProgram \cr
\noalign{\vspace{2mm}}
Studijní obor:&\StudijniObor \cr
}}}}

\vfill

% Zde doplňte rok
Praha \RokOdevzdani

\end{center}

\newpage

%%% Následuje vevázaný list -- kopie podepsaného "Zadání bakalářské práce".
%%% Toto zadání NENÍ součástí elektronické verze práce, nescanovat.

%%% Strana s čestným prohlášením k bakalářské práci

\openright
\hypersetup{pageanchor=true}
\pagestyle{plain}
\pagenumbering{roman}
\vglue 0pt plus 1fill

\noindent
Prohlašuji, že jsem tuto bakalářskou práci vypracoval(a) samostatně a výhradně
s~použitím citovaných pramenů, literatury a dalších odborných zdrojů.
Tato práce nebyla využita k získání jiného nebo stejného titulu.

\medskip\noindent
Beru na~vědomí, že se na moji práci vztahují práva a povinnosti vyplývající
ze zákona č. 121/2000 Sb., autorského zákona v~platném znění, zejména skutečnost,
že Univerzita Karlova má právo na~uzavření licenční smlouvy o~užití této
práce jako školního díla podle §60 odst. 1 autorského zákona.

\vspace{10mm}

\hbox{\hbox to 0.5\hsize{%
V \hbox to 6em{\dotfill} dne \hbox to 6em{\dotfill}
\hss}\hbox to 0.5\hsize{\dotfill\quad}}
\smallskip
\hbox{\hbox to 0.5\hsize{}\hbox to 0.5\hsize{\hfil Podpis autora\hfil}}

\vspace{20mm}
\newpage

%%% Poděkování

\openright

\noindent
\Podekovani

\newpage

%%% Povinná informační strana bakalářské práce

\openright

\vbox to 0.5\vsize{
\setlength\parindent{0mm}
\setlength\parskip{5mm}

Název práce:
\NazevPrace

Autor:
\AutorPrace

\TypPracoviste:
\Katedra

Vedoucí bakalářské práce:
\Vedouci, \KatedraVedouciho

Abstrakt:
\Abstrakt

Klíčová slova:
\KlicovaSlova

\vss}\nobreak\vbox to 0.49\vsize{
\setlength\parindent{0mm}
\setlength\parskip{5mm}

Title:
\NazevPraceEN

Author:
\AutorPrace

\TypPracovisteEN:
\KatedraEN

Supervisor:
\Vedouci, \KatedraVedoucihoEN

Abstract:
\AbstraktEN

Keywords:
\KlicovaSlovaEN

\vss}

\newpage

\openright
\pagestyle{plain}
\pagenumbering{arabic}
\setcounter{page}{1}


%%% Strana s automaticky generovaným obsahem bakalářské práce

\tableofcontents

%%% Jednotlivé kapitoly práce jsou pro přehlednost uloženy v samostatných souborech
\chapter*{Úvod}
\addcontentsline{toc}{chapter}{Úvod}

V dnešní době stále přibývá možností, kde se snažíme aplikovat metody umělé
inteligence pro řešení různorodých problémů. Na řadu z těchto problémů se nám
může nabízet hned několik možných řešení. Problém ale může nastat, pokud si
nejsme jistí, nebo třeba vůbec není možné přesně definovat, co vlastně by mělo
být správným řešením.

\paragraph{}
Pro tyto problémy se nám často hodí využívat metod evolučních algoritmů. Jedná
se o optimalizační algoritmy inspirovaných v přírodě -- specificky Darwinovou
evoluční teorií, které napodobováním přírodních procesů hledají dle našich
požadavků ta nejlepší řešení.

Zacházení s těmito algoritmy ale nemusí být vůbec jednoduché a podobně jako u
dalších optimalizačních metod a metod strojového učení je jejich běh ukryt pod
množstvím parametrů, které spolu souvisí často špatně předvídatelným způsobem.

\paragraph{}
Z tohoto důvodu je cílem této práce vytvořit platformu, která bude přístupná
uživatelům různých úrovní specializace, umožňující tvořit a provádět
experimenty s evolučními algoritmy. 

S tímto cílem volíme tvořit experimenty s roboty ve virtuálním prostředí, což
se na tento problém velmi dobře hodí. Uživatel díky robotům intuitivně chápe
složitost problému a každý posun v řešeném problému je interaktivně
pozorovatelný v daném prostředí ještě za doby řešení.

Cílem pro projekt je, aby uživatel dostal kontrolu nad experimenty a mohl tak
získat lepší přehled o práci s evolučními algoritmy a pochopil tak množství
parametrů a jejich vzájemných souvislostí, se kterými se můžeme při tvorbě
experimentů setkat. Projekt bude obsahovat různorodou řadu problémů, na které
bude potřeba využít vícero různých přístupů, což umožní dále rozšířit pochopení
problémů evolučních algoritmů. Nejtěžšími pak mohou být problémy, vyžadující
využití neuronových sítí, což může být pro uživatele díky tomuto projektu
jednoduchým, prvním využitím pokročilého algoritmu pro neuroevoluci (evoluční
vývoj neuronových sítí) -- NEAT.

Pro lehce pokročilého uživatele bude projekt dále nabízet možnosti nahlédnutí
do útrob projektu, prezentující jak se s evolučními algoritmy zachází v
samotném zdrojovém kódu, což mu zároveň umožní tvořit vlastní pokročilé
experimenty a upravovat a rozšiřovat připravenou databázi částí evolučních
algoritmů o vlastní.

\paragraph{}
Práce je rozdělena do čtyř hlavních kapitol. V kapitole
\ref{chapter-zakladní pojmy} si představíme a vysvětlíme základní pojmy
využívané v této práci jako jsou evoluční algoritmy, neuronové sítě a další.
Dále v kapitole \ref{chapter-specifikace} blíže popíšeme specifikaci projektu a
přesně jakých cílů tímto projektem chceme dosáhnout. V kapitole
\ref{chapter-implementace} již rámcově představíme jak je celý projekt interně
poskládaný a poskytneme tak základní náhled na to, jak knihovna pracuje uvnitř.
V poslední kapitole \ref{chapter-experimenty} na příkladech předvedeme
typy experimentů, které si uživatel bude moci sám ve výchozí verzi vyzkoušet, a
které bude mít možnost hned upravovat a testovat. Zároveň v této kapitole
ukážeme způsob, jakým tyto experimenty jdou statisticky vyhodnocovat a
rozebereme výsledky ukázkových experimentů. V závěru pak bude následovat
shrnutí výsledků práce a budou navržena možná rozšíření projektu.

%%% Fiktivní kapitola s ukázkami sazby

\chapter{Základní pojmy}

\section{Evoluční algoritmy}

% \section{Simulované fyzikální prostředí}

% \section{Roboti}


%%% Fiktivní kapitola s ukázkami citací

\chapter{Odkazy na literaturu}

Odkazy na literaturu vytváříme nejlépe pomocí příkazů
\verb|\citet|, \verb|\citep| atp.
(viz {\LaTeX}ový balíček \textsf{natbib}) a~následného použití
Bib{\TeX}u. V~matematickém textu obvykle odkazujeme stylem \uv{Jméno
autora/autorů (rok vydání)}, resp. \uv{Jméno autora/autorů [číslo
odkazu]}. V~českém/slovenském textu je potřeba se navíc vypořádat
s~nutností skloňovat jméno autora, respektive přechylovat jméno
autorky. Je potřeba mít na paměti, že standardní příkazy
\verb|\citet|, \verb|\citep|
produkují referenci se jménem autora/autorů v~prvním pádě a~jména
autorek jsou nepřechýlena.

Pokud nepoužíváme bib\TeX{}, řídíme se normou ISO 690 a zvyklostmi
oboru.

Jména časopisů lze uvádět zkráceně, ale pouze v~kodifikované podobě.

\section{Několik ukázek}

Mezi nejvíce citované statistické články patří práce Kaplana a~Meiera a~Coxe
\citep{KaplanMeier58, Cox72}. \citet{Student08} napsal článek o~t-testu.

Prof. Anděl je autorem učebnice matematické statistiky
\citep[viz][]{Andel98}. Teorii odhadu se věnuje práce
\citet{LehmannCasella98}. V~případě odkazů na specifickou informaci
(definice, důkaz, \dots) uvedenou v~knize bývá užitečné uvést
specificky číslo kapitoly, číslo věty atp. obsahující požadovanou
informaci, např. viz \citet[Věta 4.22]{Andel07} nebo \citep[viz][Věta
4.22]{Andel07}.

Mnoho článků je výsledkem spolupráce celé řady osob. Při odkazování
v~textu na článek se třemi autory obvykle při prvním výskytu uvedeme
plný seznam: \citet*{DempsterLairdRubin77} představili koncept EM
algoritmu. Respektive: Koncept EM algoritmu byl představen v~práci
Dempstera, Lairdové a~Rubina \citep*{DempsterLairdRubin77}. Při každém
dalším výskytu již používáme zkrácenou verzi:
\citet{DempsterLairdRubin77} nabízejí též několik příkladů použití EM
algoritmu. Respektive: Několik příkladů použití EM algoritmu lze
nalézt též v~práci Dempstera a~kol. \citep{DempsterLairdRubin77}.

U~článku s~více než třemi autory odkazujeme vždy zkrácenou formou:
První výsledky projektu ACCEPT jsou uvedeny v~práci Genbergové a~kol.
\citep{Genberget08}. V~textu \emph{nenapíšeme}: První výsledky
projektu ACCEPT jsou uvedeny v~práci \citet*{Genberget08}.

%%% Fiktivní kapitola s ukázkami tabulek, obrázků a kódu

\chapter{Implementace}

\section{Základní knihovna}

\section{Grafické rozhraní}

%%% Fiktivní kapitola s instrukcemi k PDF/A

\chapter{Experimenty a výsledky}

\section{Vývoj řížení robotů}

\section{Vývoj řízení a morfologie robotů}

\section{Diskuze výsledků}


\chapter*{Závěr}
Práce představuje platformu pro tvorbu a provádění experimentů s evolučními
algoritmy na simulovaných robotech. Práce splnila cíle, které pro ni byly
stanoveny. Vytvořili jsme systém přístupný uživatelům s různými úrovněmi
pochopení problematiky evolučních algoritmů a zpřístupnili tak nástroj, pomocí
kterého budou uživatelé moci jednoduše prohlubovat a aplikovat své znalosti.
Platforma nabízí řadu implementovaných operátorů evolučních algoritmů a několik
robotů různých úrovní složitosti, se kterými může uživatel okamžitě pracovat.

Při tvorbě platformy jsme dbali na to, aby zdrojový kód byl srozumitelný
a~jednotlivé části logicky rozdělené, což uživatelům umožní a zjednoduší
přístup ke zdrojovému kódu. Tím nabízíme možnost hlubšího pochopení fungování
evolučních algoritmů a jejich případného rozšiřování a aplikování nových
vlastních částí.

Experimenty provedené v této práci ukazují jak typy experimentů,
které~platforma umožňuje, tak styly jejich vyhodnocení. Pomocí experimentů jsme
na dvou robotech různých složitostí ověřili fakt, že jednoduché evoluční
algoritmy nestačí pro řešení všech problémů, a že pro určité pokročilé problémy
je třeba pokročilých evolučních algoritmů. 

Zároveň jsme demonstrovali příklady experimentů, které umožňují vývoj jak
řízení, tak morfologie robotů. Pro tento problém jsme v práci navrhli a
implementovali způsob, jakým je možné za běhu algoritmů měnit XML konfigurační
soubory robotů z knihovny \emph{MuJoCo}. Tyto experimenty předvedly zajímavé
výsledky, vytvářející roboty různých konfigurací, které byly schopné dosáhnout
lepších výsledků než ty výchozí.

Práce nabízí mnoho možností pro rozšiřování. Vedle pokročilejších evolučních
algoritmů na vývoj řízení robotů by zajímavým rozšířením mohlo být prozkoumání dalších
možností vývoje morfologie, umožňující rozsáhlejší změny v konfiguraci robotů.
Dalším by mohlo být podrobnější zkoumání dopadů různých fitness funkcí na
výsledky evolučního vývoje.

\addcontentsline{toc}{chapter}{Závěr}


%%% Seznam použité literatury
%%% Seznam použité literatury (bibliografie)
%%%
%%% Pro vytváření bibliografie používáme bibTeX. Ten zpracovává
%%% citace v textu (např. makro \cite{...}) a vyhledává k nim literaturu
%%% v souboru literatura.bib.
%%%
%%% Příkaz \bibliographystyle určuje, jakým stylem budou citovány odkazy
%%% v textu. V závorce je název zvoleného souboru .bst. Styly plainnat
%%% a unsrt jsou standardní součástí latexových distribucí. Styl czplainnat
%%% je dodáván s touto šablonou a bibTeX ho hledá v aktuálním adresáři.

\bibliographystyle{czplainnat}    %% Autor (rok) s českými spojkami
% \bibliographystyle{plainnat}    %% Autor (rok) s anglickými spojkami
% \bibliographystyle{unsrt}       %% [číslo]

\renewcommand{\bibname}{Seznam použité literatury}

%%% Vytvoření seznamu literatury. Pozor, pokud jste necitovali ani jednu
%%% položku, seznam se automaticky vynechá.

\bibliography{literatura}

%%% Kdybyste chtěli bibliografii vytvářet ručně (bez bibTeXu), lze to udělat
%%% následovně. V takovém případě se řiďte normou ISO 690 a zvyklostmi v oboru.

% \begin{thebibliography}{99}
%
% \bibitem{lamport94}
%   {\sc Lamport,} Leslie.
%   \emph{\LaTeX: A Document Preparation System}.
%   2. vydání.
%   Massachusetts: Addison Wesley, 1994.
%   ISBN 0-201-52983-1.
%
% \end{thebibliography}


%%% Obrázky v bakalářské práci
%%% (pokud jich je malé množství, obvykle není třeba seznam uvádět)

% \listoffigures

%%% Tabulky v bakalářské práci (opět nemusí být nutné uvádět)
%%% U matematických prací může být lepší přemístit seznam tabulek na začátek práce.

% \listoftables

%%% Použité zkratky v bakalářské práci (opět nemusí být nutné uvádět)
%%% U matematických prací může být lepší přemístit seznam zkratek na začátek práce.
% \chapwithtoc{Seznam použitých zkratek}

%%% Přílohy k bakalářské práci, existují-li. Každá příloha musí být alespoň jednou
%%% odkazována z vlastního textu práce. Přílohy se číslují.
%%%
%%% Do tištěné verze se spíše hodí přílohy, které lze číst a prohlížet (dodatečné
%%% tabulky a grafy, různé textové doplňky, ukázky výstupů z počítačových programů,
%%% apod.). Do elektronické verze se hodí přílohy, které budou spíše používány
%%% v elektronické podobě než čteny (zdrojové kódy programů, datové soubory,
%%% interaktivní grafy apod.). Elektronické přílohy se nahrávají do SISu a lze
%%% je také do práce vložit na CD/DVD. Povolené formáty souborů specifikuje
%%% opatření rektora č. 72/2017.
\appendix
\chapter{Přílohy}

\section{První příloha}

\openright
\end{document}
